\section{The weighted-$\ell^1$ entry step: the $\ell^1$ hybrid of the v2 theorem (package B10)}\label{bp:p2s2l1}

This part adds to the v2 theorem (Part~\ref{bp:p2s2}) an \emph{$\ell^1$ entry step}: the first Newton step solves its two
cohomological equations, takes its average and bounds the Lagrangian defect of the initial torus from \emph{weighted-$\ell^1$ data of the
initial torus}, with the sup-over-modes small-divisor constant
\[
\Gamma_\infty(\omega,\delta)=\sup_{k\ne0}\frac{e^{-2\pi\delta\abs k_1}}{2\pi\abs{k\cdot\omega}}\qquad(\Lean{KAM.smallDivisorSup}\ \omega\ \delta,\ \text{SHARP-L9}),
\]
which needs no lattice counting. Every other estimate of the first step, and every later step, is the sup-norm estimate of
Part~\ref{bp:p2s2}. The conclusion is that of the frozen G1 (Theorem~\ref{bp:p0:g1}); the added hypotheses bound explicit objects of the
initial torus, which the certificate computes anyway.

\emph{What is measured where.}
\begin{itemize}
\item In $\ell^1$ (initial torus $K$ only): the normal and tangential errors $\eta^N,\eta^L$ on $T_\rho$, the torsion $T$ on $T_{\rho-\delta_0}$,
  optionally the frames $L,N$ on $T_\rho$ (route A); the solutions $u=\mathcal R\eta^N$, $\xi^N$, $\xi^L$ of the first step and its increment
  $\Delta K$ are then known in $\ell^1$.
\item In sup: the Hamiltonian bounds (frozen and refined, (R1)), the frames and their caps, the torsion caps and $\avg T^{-1}$, the invariance
  error $E$ of every torus, every object of the tori $K_s$, $s\ge1$, and of the intermediate tori.
\end{itemize}
\emph{Why only at the entry} (Remark~\ref{bp:p2s2l1:design}): at later stages the errors are abstract functions whose only bounds come from
the step identities, and these contain the Hamiltonian composed with the new torus. With the frozen sup data of $h$ (D12) their $\ell^1$
norms need a sup$\to\ell^1$ conversion costing $\coth(\pi\sigma)^n\ge10^{3}$ at the available strip losses, against a small-divisor gain
of at most $6$--$18$. No conversion is used anywhere in this part.

\emph{Numbers} (\texttt{thresholds\_v2.py --l1}, \texttt{dominance.py --plasmalean}; design mirror, untrusted):
\begin{itemize}
\item PlasmaLean's example, a \emph{conditional comparison of scalar error tolerances at assumed data} (its proved tolerance $10^{-18.24}$;
  its declared bounds read as both sup and $\ell^1$ bounds, $\nu_K=\sigma_K$, $\nu_N=\sigma_N$; Remark~\ref{bp:p2s2l1:register}(i)): route A with the
  entry step and the $\ell^1$ frame increment gives $10^{-15.40}$ for the $\ell^1$ error norm ($+2.84$ orders) and $10^{-15.93}$ for the sup
  error at an assumed ratio $\ell^1/\text{sup}=3.6$ ($+2.31$); the sup-norm v2 gives $10^{-17.27}$ ($+0.97$). This does not establish dominance
  for PlasmaLean's torus (audit AUDIT-B10D-1, F7).
\item R1 (measured $\ell^1$ data, heuristic until enclosed): route A from $10^{-11.18}$ to $10^{-10.13}$ (fixed) and from $10^{-11.37}$ to
  $10^{-10.32}$ (iso), $+1.05$ orders; route B gains $0.1$.
\end{itemize}
Results of Parts~\ref{bp:p1}, \ref{bp:p2}, \ref{bp:p2s}, \ref{bp:p2s2sd} and~\ref{bp:p2s2} are cited by label.

%----------------------------------------------------------------------------------------------------
\subsection{Weighted-$\ell^1$ norms and the toolkit}\label{bp:p2s2l1:sec-norms}

\begin{definition}[Weighted-$\ell^1$ norms on strips]\label{bp:p2s2l1:norms}
For $r>0$ and $f:T_r\to\C$ holomorphic and $1$-periodic with Fourier coefficients $\hat f_k$ on $\R^n$, put
$\norm f_{1,r}:=\sum_k\abs{\hat f_k}e^{2\pi r\abs k_1}\in[0,\infty]$ (the repository's \Lean{Fourier.weightedNorm} at the coefficient radius $2\pi r$ of
\Lean{torusCoeff}). For $f:T_r\to\C^p$, $\norm f_{1,r}:=\max_i\norm{f_i}_{1,r}$; for a matrix function
$M:T_r\to\C^{p\times q}$, $\norm M_{1,r}:=\max_i\sum_j\norm{M_{ij}}_{1,r}$ (row sums). A bound $\norm f_{1,r}\le A$ always includes the
summability of the coefficient families and $A\ge0$ (bound predicates, as \Lean{WBound}; no junk values).
\end{definition}

\begin{lemma}[The $\ell^1$ toolkit]\label{bp:p2s2l1:toolkit}
Let $r>0$, and $f,g,M$ holomorphic and $1$-periodic on $T_r$ with finite norms.
\begin{enumerate}[label=(\alph*)]
\item \emph{(sup bound)} $\abs{f(\theta)}\le\norm f_{1,r}$ and $\abs{M(\theta)}\le\norm M_{1,r}$ for $\theta\in T_r$ (max norm and its operator
  norm, the maximal row sum);
\item \emph{(products)} $\norm{Mg}_{1,r}\le\norm M_{1,r}\norm g_{1,r}$ and $\norm{fg}_{1,r}\le\norm f_{1,r}\norm g_{1,r}$; a constant
  vector $c$ has $\norm c_{1,r}=\abs c$; $\norm f_{1,r'}\le\norm f_{1,r}$ for $r'\le r$; $\norm{f-\avg f}_{1,r}\le\norm f_{1,r}$;
\item \emph{(row-sum Cauchy)} for $0<\delta<r$ and $f:T_r\to\C^p$: on $T_{r-\delta}$, $\abs{Df(\theta)}\le\norm f_{1,r}/(e\delta)$ and
  $\abs{(Df(\theta))^\top}\le p\,\norm f_{1,r}/(e\delta)$; for $p=n$, $\abs{Df(\theta)-(Df(\theta))^\top}\le n\norm f_{1,r}/(e\delta)$.
\end{enumerate}
\source{(a): Lean \Lean{KAM.norm\_le\_weightedNorm\_of\_summableWeighted} and its vector and matrix forms \Lean{KAM.norm\_le\_of\_l1\{,\_pi,\_mat\}}
(SHARP-L9). (b): \Lean{KAM.Fourier.weightedNorm\_convolution\_le} with \Lean{KAM.Fourier.convolution}. (c): \Lean{KAM.Fourier.weightedNorm\_partialCoeff\_le}
(coefficientwise), summed over $j$ as in\\ \Lean{KAM.Weighted.cauchy\_theta}. Ours (assembly).}
\end{lemma}

\begin{proof}
\proofstep{Sup bound}\label{step:p2s2l1:tk-sup}
A holomorphic periodic $f$ with summable weighted coefficients equals its Fourier series on $T_r$ (\Lean{eq\_fourierSum\_of\_summableWeighted}),
and $\abs{e^{2\pi ik\cdot\theta}}\le e^{2\pi r\abs k_1}$ there. For matrices apply this to each entry and sum a row.

\proofstep{Products}\label{step:p2s2l1:tk-prod}
On $\R^n$ both factors equal their absolutely convergent Fourier series, so the coefficients of the product are the convolutions of the
coefficients; $e^{2\pi r\abs{k}_1}\le e^{2\pi r\abs l_1}e^{2\pi r\abs{k-l}_1}$ gives submultiplicativity entry by entry, and
$\norm{(Mg)_i}\le\sum_j\norm{M_{ij}}\norm{g_j}\le(\sum_j\norm{M_{ij}})\max_j\norm{g_j}$. Removing the zero mode removes one nonnegative
term.

\proofstep{Row-sum Cauchy}\label{step:p2s2l1:tk-cauchy}
$\widehat{\partial_jf_i}_k=2\pi ik_j\hat f_{i,k}$ and $\sum_j2\pi\abs{k_j}e^{-2\pi\delta\abs k_1}=2\pi\abs k_1e^{-2\pi\delta\abs k_1}\le(e\delta)^{-1}$,
so $\sum_j\norm{\partial_jf_i}_{1,r-\delta}\le\norm{f_i}_{1,r}/(e\delta)$; (a) gives the first bound. For the transpose,
$\sum_i\abs{\partial_jf_i}\le\sum_i\norm{f_i}/(e\delta)$. For the antisymmetric part, row $i$ is bounded by
$\sum_j\norm{\partial_jf_i}+\sum_{j\ne i}\norm{\partial_if_j}\le\norm{f_i}/(e\delta)+(n-1)\max_j\norm{f_j}/(e\delta)$.
\end{proof}

\begin{lemma}[Conversion and composition in $\ell^1$: the prices]\label{bp:p2s2l1:convert}
\begin{enumerate}[label=(\alph*)]
\item If $f$ is holomorphic, periodic and $\abs f\le A$ on $T_r$, then $\norm f_{1,r-\sigma}\le\coth(\pi\sigma)^nA$ for $0<\sigma<r$, and
  $\coth(\pi\sigma)\ge(\pi\sigma)^{-1}$.
\item If $\xi:T_{r'}\to\C^n$ is periodic with $r'+\norm\xi_{1,r'}<r$, then $\norm{f\circ(\mathrm{id}+\xi)}_{1,r'}\le\norm f_{1,r'+\norm\xi_{1,r'}}$
  and $\norm{f\circ(\mathrm{id}+\xi)-f}_{1,r'}\le\norm\xi_{1,r'}\norm f_{1,r}/\bigl(e(r-r'-\norm\xi_{1,r'})\bigr)$.
\end{enumerate}
Neither is used by the theorems of this part; they price the alternatives of Remark~\ref{bp:p2s2l1:design}.
\source{(a): Lemma~\ref{bp:p1:summation} (Lean \Lean{KAM.Weighted.wbound\_of\_sup\_bound}); PLAN \S3. (b): $e^{2\pi ik\cdot(\theta+\xi)}=e^{2\pi ik\cdot\theta}\exp(2\pi ik\cdot\xi)$
and $\norm{\exp(g)-1}_{1}\le\norm g_1e^{\norm g_1}$ in the Wiener algebra (PLAN \S3, composition).}
\end{lemma}

%----------------------------------------------------------------------------------------------------
\subsection{The $\ell^1$ solves}\label{bp:p2s2l1:sec-solve}

\begin{lemma}[Solutions, averages and the ray]\label{bp:p2s2l1:solve}
Let $\omega\in\Dio_{\gamma,\tau}$, $0<\delta<r$, and $g:T_r\to\C^p$ holomorphic, periodic, with $\avg g=0$ and $\norm g_{1,r}\le A$. Then
\begin{enumerate}[label=(\alph*)]
\item $\norm{\mathcal Rg}_{1,r-\delta}\le\Gamma_\infty(\omega,\delta)A$, hence $\abs{\mathcal Rg}\le\Gamma_\infty(\omega,\delta)A$ on $T_{r-\delta}$;
\item $\abs{\mathcal Rg(\theta)}\le\Gamma_\infty(\omega,r)A$ for $\theta\in\R^n$ (the full bite: radius $0$); hence for a continuous periodic matrix function
  $Y$ with $\abs Y\le y$ on $\R^n$, $\abs{\avg{Y\mathcal Rg}}\le y\,\Gamma_\infty(\omega,r)A$;
\item $\Gamma_\infty(s\omega,\delta)=s^{-1}\Gamma_\infty(\omega,\delta)\le\Gamma_\infty(\omega,\delta)$ for $s\ge1$: a certificate for $\omega^*$ serves the whole
  ray $\{s\omega^*:s\ge1\}$.
\end{enumerate}
\source{(a): Lean \Lean{KAM.cohomological\_l1} (vector clauses; SHARP-L9). (b): the coefficient estimate
\Lean{KAM.summableWeighted\_solve\_sup} at radius $0$ (it needs no $\delta<r$) and\\ \Lean{KAM.hasSum\_torusCoeff\_mul\_character} on $\R^n$;
Corollary~\ref{bp:p2:averages}(a). (c): ours.}
\end{lemma}

\begin{proof}
\proofstep{Coefficients}\label{step:p2s2l1:so-coef}
$\widehat{\mathcal Rg}_k=\hat g_k/(2\pi ik\cdot\omega)$ for $k\ne0$, so, term by term,
\[
\abs{\widehat{\mathcal Rg}_k}e^{2\pi(r-\delta)\abs k_1}\le\Gamma_\infty(\omega,\delta)\abs{\hat g_k}e^{2\pi r\abs k_1};
\]
sum, and apply Lemma~\ref{bp:p2s2l1:toolkit}(a). At $\delta=r$ the same computation bounds $\norm{\mathcal Rg}_{1,0}$, which bounds
$\mathcal Rg$ on $\R^n$ where $\mathcal Rg$ is the sum of its absolutely convergent Fourier series.

\proofstep{The ray}\label{step:p2s2l1:so-ray}
$\abs{k\cdot s\omega}=s\abs{k\cdot\omega}$ term by term.
\end{proof}

\begin{lemma}[The Lagrangian defect from $u=\mathcal R\eta^N$]\label{bp:p2s2l1:defect}
Under the standing hypotheses of Part~\ref{bp:p2} (in particular $h$ periodic and real, $K$ primary), let $K$ be holomorphic on $T_r$ with
$\det G\ne0$ and $\norm{\eta^N}_{1,r}\le\varepsilon^1_N$, and put $u:=\mathcal R\eta^N$. Then $\Omega_K=-(Du-(Du)^\top)$ on $T_{r-\delta}$ for every
$0<\delta<r$, and for $0<\delta'<r-\delta$, on $T_{r-\delta-\delta'}$,
\[
\abs{\Omega_K}\le\frac{n}{e\,\delta'}\,\Gamma_\infty(\omega,\delta)\,\varepsilon^1_N .
\]
\source{Lemma~\ref{bp:p2s2:omega} (FH24 Lemma~3.1) with the sign of Lean \Lean{KAM.Sharp2.fderiv\_lagrangeMat\_eq\_etaN} (SHARP2-B2; audit
AUDIT-SHARP-V2-iso-1, F8); ours.}
\end{lemma}

\begin{proof}
\proofstep{The identity}\label{step:p2s2l1:de-id}
$\avg{\eta^N}=0$ (Lemma~\ref{bp:p2:exactness}), so $\partial_\omega u=\eta^N$ and $\partial_\omega(Du-(Du)^\top)=D\eta^N-(D\eta^N)^\top=-\partial_\omega\Omega_K$. Both
$\Omega_K$ (Lemma~\ref{bp:p2:exactness}) and $Du-(Du)^\top$ (derivatives of periodic functions) have zero average; the uniqueness clause of
Lemma~\ref{bp:p1:cohomological} gives the identity.

\proofstep{The bound}\label{step:p2s2l1:de-b}
Lemma~\ref{bp:p2s2l1:solve}(a) gives $\norm u_{1,r-\delta}\le\Gamma_\infty(\omega,\delta)\varepsilon^1_N$, and Lemma~\ref{bp:p2s2l1:toolkit}(c), last clause, with
the bite $\delta'$.
\end{proof}

For $n\ge2$ the factor $n/e$ is below the factor $2(n-1)$ of Lemma~\ref{bp:p2s2:omega}, and no second bite constant is needed for the
derivative ($\delta'$ is a Cauchy loss).

%----------------------------------------------------------------------------------------------------
\subsection{Where $\ell^1$ pays: the design decision}\label{bp:p2s2l1:sec-design}

\begin{remark}[Decision and the price of each conversion]\label{bp:p2s2l1:design}
Write $\varrho:=\coth(\pi\sigma)^n$ for the conversion of Lemma~\ref{bp:p2s2l1:convert}(a). For $n=2$: R1 ($\delta_0=0.00238$) $\varrho=1.8\cdot10^4$
at $\sigma=\delta_0$, $2.0\cdot10^3$ at $3\delta_0$, $5.0\cdot10^2$ at $\rho$; PlasmaLean ($\delta_0=0.00333$) $9.1\cdot10^3$, $1.0\cdot10^3$, $2.5\cdot10^2$. The
small-divisor gain is $c_{\mathrm{freq}}/c_\infty=5.9$ (R1, $0.0030/0.00051$) and $6.0$ (PlasmaLean, $0.0075/0.00126$), at most $18$ against the
generic $\ell^1$ constant. The decision for each object that feeds a solve:
\begin{center}\small
\begin{tabular}{p{3.3cm}p{2.5cm}p{3.6cm}p{4.0cm}}
object & feeds & norm used & conversion\\\hline
$\eta^N_0,\eta^L_0$ (stage $0$) & $u$, $g^L$ & $\ell^1$, certificate & none\\
$T_0$ (stage $0$) & $g^L=\eta^L-T\xi^N$ & $\ell^1$, certificate & none\\
$\xi^N_0$, $u$, $\xi^L_0$ & $g^L$, $\Delta K_0$ & $\ell^1$, from the solves & none\\
$\eta^N_s,\eta^L_s$, $s\ge1$ & $u_s$, $g^L_s$ & sup & none ($\ell^1$ would cost $\varrho\ge10^3$)\\
$\Delta^2X$, $X_N\xi^N$, $\tilde K\circ\varphi$ & new error & sup & none ($\ell^1$ would cost $\varrho$, or $h$ in a Fourier--Taylor norm)\\
$L,N,B$, $T_s$, $\avg T^{-1}$ & caps & sup ($L,N$ also $\ell^1$ at stage $0$, route A) & none
\end{tabular}
\end{center}
\begin{enumerate}[label=(\roman*)]
\item \emph{Stage $s\ge1$ in $\ell^1$ does not pay.} In route A the new error is $\bar E=PS^{-1}E_{\mathrm{lin}}+\Delta^2X$; on R1 with the levers of
  Part~\ref{bp:p2s2} the Taylor term is $16$, $23$, $33$, $52$, $68\,\%$ of $C_E$ at stages $0$--$4$ (\texttt{stageA}). With sup data on $h$ its $\ell^1$ norm costs
  $\varrho\ge2\cdot10^3$ (at the largest available loss $3\delta_s$), so even with a free $\ell^1$ bound on $PS^{-1}E_{\mathrm{lin}}$ the $\ell^1$ error fed to stage $s+1$
  is at least $0.16\varrho>300$ times its sup bound, while the next solve gains at most $18$; the per-ingredient minimum would never select it. The
  linear part needs in addition $\norm{T_s}_1$, whose increments contain $DX_h\circ K_{s}-DX_h\circ K_{s-1}$ (again a conversion). Route B's new errors contain $\Delta^2X$, $X_N\xi^N$ (through $C_{\tilde E}$)
  and the composition $\tilde E\circ\varphi$: the same conclusion, and Lemma~\ref{bp:p2s2l1:convert}(b) would add the Wiener exponential to Lean.
  The alternative, $h$ in a Fourier--Taylor norm (PLAN \S3), is a new hypothesis outside the frozen sup data (D12), a new certificate primitive
  and a composition layer of $\approx1.5$k Lean lines; it would add at most the stage-$1$ share of the gain (a quarter of the log-gain of the
  quadratic condition, nothing for the cap sums, which stage $0$ dominates).
\item \emph{Stage $0$ in $\ell^1$ pays, with no conversion.} The certificate computes $E$, the frames and $T$ as Fourier models; $\eta=P^\top\Omega_0E$
  and $T\xi^N$ are products of such models (SHARP-CERTSIDE \S2.7). The binding constraints after the levers of Part~\ref{bp:p2s2} are the frame
  perturbation cap sums (route A), which stage $0$ dominates, and the tail decay (route B), where stage $0$ carries half of the log-weight.
\item \emph{Mixed products never occur in an $\ell^1$ estimate.} The only products estimated in $\ell^1$ are $T\xi^N$ ($\ell^1\times\ell^1$, with the
  certified $\norm T_{1,\rho-\delta_0}$) and, for the route-A frame increment, $L\xi^L$, $N\xi^N$ ($\ell^1\times\ell^1$, with $\norm L_{1,\rho},\norm N_{1,\rho}$).
  Sup bounds are used only in sup estimates; Lemma~\ref{bp:p2s2l1:toolkit}(a) turns $\ell^1$ bounds into sup bounds where a sup estimate needs
  them.
\item \emph{Each ingredient takes the smaller bound.} Both the sup constant (Definition~\ref{bp:p2s2sd:const}, at the same bites) and the $\ell^1$
  route bound the same objects, so the entry constants are minima (Definitions~\ref{bp:p2s2l1:constA}, \ref{bp:p2s2l1:constB}), and the entry
  is never worse than the sup stage (Lemma~\ref{bp:p2s2l1:never-worse}). On R1 the full-strip average is the ingredient where sup wins
  ($c(\rho)=0.0067$ against $a_N c_\infty(\rho)=4.1\cdot0.0020$).
\end{enumerate}
\end{remark}

%----------------------------------------------------------------------------------------------------
\subsection{The entry data}\label{bp:p2s2l1:sec-data}

\begin{assumption}[$\ell^1$ entry data]\label{bp:p2s2l1:data}
In addition to Assumption~\ref{bp:p2s2:standing}, with the frozen bound $\abs E\le e$ on $T_\rho$, the stage-$0$ bite $\delta_0$ and the frame caps
$\sigma_{L^\top},\sigma_{N^\top}$ of \ref{bp:p2s2:R2}, the following numbers are given, positive (well-formedness), with the hypotheses stated:
\begin{enumerate}[label=(L\arabic*)]
\item\label{bp:p2s2l1:L1} $a_N,a_L>0$ with $\norm{\eta^N}_{1,\rho}\le a_N\sigma_{L^\top}e$ and $\norm{\eta^L}_{1,\rho}\le a_L\sigma_{N^\top}e$;
\item\label{bp:p2s2l1:L2} $t_1$ with $\norm T_{1,\rho-\delta_0}\le t_1$, and $t_{\R}\le t_1$ with $\abs T\le t_{\R}$ on $\R^n$ ($T$ the frozen torsion of $K$ at $\omega$; $t_1$ need not
  dominate a derived sup cap $t$);
\item\label{bp:p2s2l1:L3} $G_0,G_0',G_\rho>0$ with $\Gamma_\infty(\omega,\delta_0)\le G_0$, $\Gamma_\infty(\omega,\delta_0/2)\le G_0'$, $\Gamma_\infty(\omega,\rho)\le G_\rho$ (for the actual
  $\omega$; on the ray of $\omega^*$ by Lemma~\ref{bp:p2s2l1:solve}(c); certified by \Lean{KAM.smallDivisorSup\_snoc\_le});
\item\label{bp:p2s2l1:L4} \emph{(optional, route A)} $\nu_K,\nu_N>0$ with $\norm L_{1,\rho}\le\nu_K$, $\norm N_{1,\rho}\le\nu_N$;
\item\label{bp:p2s2l1:L5} the sup small-divisor constants of \ref{bp:p2s2:R4} at the stage-$0$ bites $\delta_0,\delta_0/2,\rho$ and at every later prefix bite,
  each an $\mathrm{SDConst}$ premise for the actual $\omega$ (or its ray) at that bite (audit AUDIT-SHARP-V2-fixed-1, F3, F9).
\end{enumerate}
Put $c^\infty_0:=G_0\bar\gamma\delta_0^\varkappa$, $c^{\infty\prime}_0:=G_0'\bar\gamma(\delta_0/2)^\varkappa$, $c^\infty_\rho:=G_\rho\bar\gamma\rho^\varkappa$. Each hypothesis bounds an explicit
object of $(h,K,\omega)$ or a linear solver; none is conclusion-shaped. Nothing here is implied by the frozen hypotheses without a conversion;
the entry is used only when its data are certified.
\end{assumption}

For R1 (\texttt{l1\_entry\_r1.py}, sampled, heuristic): $\norm E_{1,\rho}/\sup\abs E=4.28$, $a_N=3.92$, $a_L=0.26$, $\norm T_{1,\rho-\delta_0}=3.55$
($\sup=3.31$), $\norm L_{1,\rho}=1.36$, $\norm N_{1,\rho}=1.42$ ($1.04\times$ the sup norms); $c^\infty_0=0.00051$, $c^{\infty\prime}_0=0.00049$, $c^\infty_\rho=0.0020$.

%----------------------------------------------------------------------------------------------------
\subsection{Route A: the entry step}\label{bp:p2s2l1:sec-A}

Route A at stage $0$ is the additive step (Proposition~\ref{bp:p2s:step}) with the constants of Definition~\ref{bp:p2s2:constantsA} (mirror
\texttt{stageA}); write $u_0:=e/(\gamma\delta_0^\varkappa)$, $g_0:=\bar\gamma\delta_0^\varkappa$, and $c,\bar c$ for its sup constants at $\delta_0$ and $\rho$.

\begin{definition}[Entry constants, route A]\label{bp:p2s2l1:constA}
The constants of Definition~\ref{bp:p2s2:constantsA} at stage $0$ with the following replacements, in this order:
\begin{align*}
&c_{\mathrm{avg}}:=\min\{\bar c,\ a_Nc^\infty_\rho\}(\delta_0/\rho)^\varkappa,\quad c_N:=\min\{c,\ a_Nc^\infty_0\},\quad c^1_N:=a_Nc^\infty_0,\\
&C_S:=\sigma_T(\sigma_{N^\top}g_0+\min\{t,t_{\R}\}\,c_{\mathrm{avg}}\sigma_{L^\top}),\quad C_{\xi N}:=C_S+c_N\sigma_{L^\top},\quad C^1_{\xi N}:=C_S+c^1_N\sigma_{L^\top},\\
&C_g:=t\,C_{\xi N}+\sigma_{N^\top}g_0,\quad C^1_{\xi L}:=c^\infty_0\bigl(t_1C^1_{\xi N}+a_L\sigma_{N^\top}g_0\bigr),\quad C_{\xi L}:=\min\{c\,C_g,\ C^1_{\xi L}\},\\
&C_\xi:=\max\{C_{\xi L},g_0C_{\xi N}\},\quad C_{\Delta K}:=\sigma_KC_{\xi L}+\sigma_Ng_0C_{\xi N},\\
&C^K_\Omega:=\min\{C_\Omega,\ \tfrac ne\,a_Nc^\infty_0\sigma_{L^\top}\},\quad C_{\mathrm{sym}}:=\max\{1,\sigma_B^2\}C^K_\Omega,
\end{align*}
and in $C_{\Delta T}$ the term $2\sigma_B^3wC_\Omega g_0$ (the parts $\mathcal T_c$ of $K$ and $\bar K$) becomes $\sigma_B^3w(C^K_\Omega+C_\Omega)g_0$. With \ref{bp:p2s2l1:L4}
(the $\ell^1$ frame increment, used only when its certificate is given) put $C^1_{\Delta K}:=\nu_KC^1_{\xi L}+\nu_Ng_0C^1_{\xi N}$ and the frame
increments $d_L:=\min\{nC_{\Delta K},C^1_{\Delta K}/e\}$, $d_{L^\top}:=\min\{2nC_{\Delta K},2nC^1_{\Delta K}/e\}$; without it $d_L:=nC_{\Delta K}$,
$d_{L^\top}:=2nC_{\Delta K}$. They replace $nC_{\Delta K}$ and $2nC_{\Delta K}$ in the frame increment: $r_B:=\sigma_B^2(\sigma_{L^\top}d_L+\sigma_Kd_{L^\top})/(1-\vartheta)$
(for $C_{\bar B}C_{\Delta K}$), $r_N:=\sigma_Bd_L+\sigma_Kr_B$, $r_{N^\top}:=r_B\sigma_{L^\top}+\sigma_Bd_{L^\top}$, the $\mathcal T_c$-change $\sigma_B^3w(\sigma_Kd_{L^\top}+\sigma_{L^\top}d_L)$,
the Gram condition $\sigma_B(\sigma_{L^\top}d_L+\sigma_Kd_{L^\top})$, and the rates $d_L/\eta$ (DK) and $d_{L^\top}/\eta$ ($L^\top$). These are products, with no
division by $C_{\Delta K}$, and an enabled frame option has $\nu_K,\nu_N>0$ (audit AUDIT-B10D-1, F1). The mirror is
\texttt{stageA(\dots, l1=\dots)}; for the isoenergetic step \texttt{iso=True} adds $\Theta C_S$ inside $C^1_{\xi L}$ and $C^K_\Omega$ replaces $C_\Omega$ in
the frequency terms (track-I part \texttt{p2s2-iso-l1hybrid}).
\end{definition}

\begin{lemma}[The $\ell^1$ frame increment]\label{bp:p2s2l1:frame}
Under \ref{bp:p2s2l1:L4}, with $\Delta K=L\xi^L+N\xi^N$ on $T_{\rho-2\delta_0}$:
$\norm{\Delta K}_{1,\rho-2\delta_0}\le C^1_{\Delta K}e/(\gamma^2\delta_0^{2\varkappa})$, and on $T_{\rho-(2+\eta)\delta_0}$
\[
\abs{D\bar K-DK}\le\frac{\norm{\Delta K}_{1,\rho-2\delta_0}}{e\eta\delta_0},\qquad\abs{(D\bar K-DK)^\top}\le\frac{2n\norm{\Delta K}_{1,\rho-2\delta_0}}{e\eta\delta_0}.
\]
\source{Lemma~\ref{bp:p2s2l1:toolkit}(b), (c); ours. Route A only: route B's increment contains the composition $\tilde K\circ\varphi$.}
\end{lemma}

\begin{proposition}[The entry step, route A]\label{bp:p2s2l1:entryA}
Let $K$ be $(\rho,\beta,e)$-admissible with the refined data and the entry data \ref{bp:p2s2l1:L1}--\ref{bp:p2s2l1:L5}. If the step conditions of
Definition~\ref{bp:p2s2:constantsA} hold at $x_0=e/(\gamma^4\delta_0^{4\varkappa})$ with the entry constants, then the conclusions (a)--(f) of
Proposition~\ref{bp:p2s:step} hold for $\bar K=K+\Delta K$ at stage $0$ with the entry constants (and those of the caps $L^\top,N,N^\top$, $\abs T$ of
Part~\ref{bp:p2s2}): $\abs{\bar E}\le C_Ee^2/(\gamma^4\delta_0^{4\varkappa})$ on $T_{\rho-2\delta_0}$ and the increments of every cap on $T_{\rho-3\delta_0}$.
\source{HL19 Lemma~5.2; Proposition~\ref{bp:p2s:step}; Definition~\ref{bp:p2s2:constantsA}; ours (the entry).}
\end{proposition}

\begin{proof}
\proofstep{The normal solve}\label{step:p2s2l1:ea-n}
By \ref{bp:p2s2l1:L1} and Lemma~\ref{bp:p2s2l1:solve}(a), $\norm u_{1,\rho-\delta_0}\le G_0a_N\sigma_{L^\top}e\le c^1_N\sigma_{L^\top}u_0$; Definition~\ref{bp:p2s2sd:const}(a)
gives $\abs u\le c\sigma_{L^\top}u_0$; so $\abs u\le c_N\sigma_{L^\top}u_0$ on $T_{\rho-\delta_0}$.

\proofstep{The average}\label{step:p2s2l1:ea-avg}
$\abs{\avg{Tu}}\le\min\{t,t_{\R}\}\min\{\bar c\,u_0(\delta_0/\rho)^\varkappa,\ G_\rho a_N e\}\sigma_{L^\top}$ ($t$ and $t_{\R}$ both bound $\abs T$ on $\R^n$) by Lemmas~\ref{bp:p2s2sd:avg} and~\ref{bp:p2s2l1:solve}(b), and
$\abs{\avg{\eta^L}}\le\sigma_{N^\top}e$; hence $\abs{\xi^N_0}\le C_Su_0$, $\abs{\xi^N}\le C_{\xi N}u_0$ and $\norm{\xi^N}_{1,\rho-\delta_0}\le C^1_{\xi N}u_0$
(a constant has $\ell^1$ norm its modulus, Lemma~\ref{bp:p2s2l1:toolkit}(b)).

\proofstep{The tangential solve}\label{step:p2s2l1:ea-l}
$g^L=T\xi^N-\eta^L$ has $\avg{g^L}=0$ and, by Lemma~\ref{bp:p2s2l1:toolkit}(b) and \ref{bp:p2s2l1:L2},
$\norm{g^L}_{1,\rho-\delta_0}\le t_1C^1_{\xi N}u_0+a_L\sigma_{N^\top}e$; Lemma~\ref{bp:p2s2l1:solve}(a) gives
$\norm{\xi^L}_{1,\rho-2\delta_0}\le C^1_{\xi L}e/(\gamma^2\delta_0^{2\varkappa})$, and the sup route gives $c\,C_g$; the sup bound $\abs{g^L}\le C_gu_0$ is unchanged.

\proofstep{The defect}\label{step:p2s2l1:ea-om}
Lemma~\ref{bp:p2s2l1:defect} with $\delta=\delta'=\delta_0$ gives, on $T_{\rho-2\delta_0}$,
$\abs{\Omega_K}\le\frac n{e\delta_0}G_0a_N\sigma_{L^\top}e\le\frac ne a_Nc^\infty_0\sigma_{L^\top}e/(\gamma\delta_0^{\varkappa+1})$, against $C_\Omega e/(\gamma\delta_0^{\varkappa+1})$ from the sup route; $E_{\mathrm{sym}}$ follows as in Lemma~\ref{bp:p2s:geometry}.

\proofstep{The rest}\label{step:p2s2l1:ea-rest}
Lemma~\ref{bp:p2s:correction} (its last two lines), Lemma~\ref{bp:p2s:newton-error} and Lemma~\ref{bp:p2s:perturbation} use only the bounds of the previous
steps, the sup data and the Taylor bound, so they hold with the entry constants; $\mathcal T_c$ of $\bar K$ is bounded through its own (sup) defect,
as in Part~\ref{bp:p2s2}. With \ref{bp:p2s2l1:L4}, Lemma~\ref{bp:p2s2l1:frame} replaces the Cauchy step on $\bar K-K$ in Lemma~\ref{bp:p2s:perturbation}.
\end{proof}

%----------------------------------------------------------------------------------------------------
\subsection{Route B: the entry step}\label{bp:p2s2l1:sec-B}

Route B at stage $0$ is Proposition~\ref{bp:p2s2:step} with the weighted error $\mu_0=pe$, $p:=\max\{\sigma_{N^\top},\sigma_{L^\top}/(\gamma\delta_0^\varkappa)\}$ (and $1$ in the iso
form), and the clipped weight $\lambda_0:=\sigma_{N^\top}/p\le1$ (audit AUDIT-SHARP-V2-iso-1, F4). Put $r_N:=a_N\sigma_{L^\top}/(\gamma\delta_0^\varkappa p)\le a_N$ and
$r_L:=a_L\lambda_0$, so that $\norm{\eta^N}_{1,\rho}\le r_N\gamma\delta_0^\varkappa\mu_0$ and $\norm{\eta^L}_{1,\rho}\le r_L\mu_0$.

\begin{definition}[Entry constants, route B]\label{bp:p2s2l1:constB}
The constants of Definition~\ref{bp:p2s2:constantsB} at stage $0$ (with the stage-$0$ replacements of $C_E$, $C_{DE}$ by $e/\mu_0$, $ne/\mu_0$), with
\begin{align*}
&\overline{\mathrm{avg}}:=\min\{\bar c,\ r_Nc^\infty_\rho\}(\delta_0/\rho)^\varkappa\ \text{in }C_0,\quad
C_{\xi N}:=C_0+\min\{c,r_Nc^\infty_0\},\quad C^1_{\xi N}:=C_0+r_Nc^\infty_0,\\
&c^{(1)}_N:=\min\{c^{(1)},\ 2^{\varkappa+1}r_Nc^{\infty\prime}_0\}\ \text{in }C_{D\tilde L},C_{D\tilde L^\top},\qquad
C_{\xi L}:=\min\bigl\{c\,C_{gL},\ c^\infty_0\bigl(r_L+t_1C^1_{\xi N}\bigr)\bigr\},\\
&C^K_\Omega:=\min\{C_\Omega,\ 2^{\varkappa+1}\tfrac ne\,r_Nc^{\infty\prime}_0\},
\end{align*}
where $C^K_\Omega$ (the defect of $K$) replaces $C_\Omega$ in the condition $\max\{1,\sigma_B^2\}C_\Omega g X\le q$, in the term $\sigma_B^2C_\Omega g^2$ of $Q_L$, in one of the
two terms $\sigma_B^3wC_\Omega$ of $C_{\Delta T}$ and in the derived-torsion condition; $C_\Omega$ (the defects of $\tilde K$ and $\bar K$, sup data) stays in the
condition $C_\Omega C_{gL}X\le\vartheta_\Omega$, in the other $\mathcal T_c$ term and in the isoenergetic terms. In the isoenergetic form $\Theta C_0$ is added inside the
$\ell^1$ branch of $C_{\xi L}$, the cubic term of $\bar Q_N$ keeps the sup value $cC_{gL}$ in its denominator, and the condition $cC_{gL}X\le\vartheta_\xi$ is
added. The mirror is \texttt{stageB(\dots, l1=\dots)}.
\end{definition}

\begin{proposition}[The entry step, route B]\label{bp:p2s2l1:entryB}
Under the hypotheses of Proposition~\ref{bp:p2s2:step} at stage $0$ and the entry data \ref{bp:p2s2l1:L1}--\ref{bp:p2s2l1:L3}, \ref{bp:p2s2l1:L5}, the
conclusions of Proposition~\ref{bp:p2s2:step} hold with the entry constants, provided the step conditions hold with them. The numbers
$p,\lambda_0,r_N,r_L$ may be replaced by any $\mu_0>0$, $0<\lambda\le1$ and upper bounds with $\sigma_{N^\top}e\le\lambda\mu_0$,
$\sigma_{L^\top}e\le\gamma\delta_0^\varkappa\mu_0$, $e\le\mu_0/p$, $\norm{\eta^N}_{1,\rho}\le r_N\gamma\delta_0^\varkappa\mu_0$, $\norm{\eta^L}_{1,\rho}\le r_L\mu_0$ (Lean
\Lean{entry\_step\_B}); in the theorem uniform in $\gamma\le\bar\gamma$ one takes $\lambda_0=\sigma_{N^\top}\bar\gamma\delta_0^\varkappa/\sigma_{L^\top}$ under the regime
$\sigma_{N^\top}\bar\gamma\delta_0^\varkappa\le\sigma_{L^\top}$, so that $\mu_0=\sigma_{L^\top}e/(\gamma\delta_0^\varkappa)$ and $r_N=a_N$.
\source{FH24 Lemma~3.6; Proposition~\ref{bp:p2s2:step}; ours (the entry).}
\end{proposition}

\begin{proof}
\proofstep{The four $\ell^1$ ingredients}\label{step:p2s2l1:eb-in}
As Steps~\ref{step:p2s2l1:ea-n}--\ref{step:p2s2l1:ea-om}: $u$ on $T_{\rho-\delta_0}$ (bite $\delta_0$); the average at real points (bite $\rho$); $D(Nu)$ from
$\abs{Nu}\le\sigma_NG_0'\norm{\eta^N}_{1,\rho}$ on $T_{\rho-\delta_0/2}$ and the sup Cauchy estimate with the bite $\delta_0/2$ (Lemma~\ref{bp:p1:cauchy-corollaries}),
which is the structure of Corollary~\ref{bp:p2s2sd:rc} with $s=\frac12$; $\xi^L=-\mathcal R(\eta^L-T\xi^N)$ on $T_{\rho-2\delta_0}$; the defect of $K$ on $T_{\rho-\delta_0}$ by
Lemma~\ref{bp:p2s2l1:defect} with $\delta=\delta'=\delta_0/2$.

\proofstep{The rest}\label{step:p2s2l1:eb-rest}
Steps~\ref{step:p2s2:st-int}--\ref{step:p2s2:st-fr} use these bounds, the sup data and the sup lemmas for $\tilde K$, $\bar K$, unchanged.
\end{proof}

\begin{lemma}[Never worse than the sup stage]\label{bp:p2s2l1:never-worse}
Assume the entry data of Assumption~\ref{bp:p2s2l1:data} are certified. Every final stage-$0$ constant of Definitions~\ref{bp:p2s2l1:constA}
and~\ref{bp:p2s2l1:constB} (the rates, $C_E$, and the coefficients of the step conditions) is at most the corresponding sup constant, every added step
condition is a step condition of the sup stage, and the stage constants are nondecreasing in the replaced quantities. Hence the conditions
(J1)--(J3) with the entry at stage $0$ are implied by those of the sup scheme. The auxiliary $\ell^1$ constants ($C^1_{\xi N}$, $C^1_{\xi L}$, $C^1_{\Delta K}$)
need not be below their sup counterparts, and the sup data do not imply the entry data on the same strip: a bounded holomorphic function on
$T_\rho$ need not have a summable weighted series at $\rho$ (audit AUDIT-B10D-1, F6). Never-worse is therefore a statement for certified entry data,
and the final theorem keeps the sup branches and takes the minimum by cases (Corollary~\ref{bp:p2s2l1:frozen}).
\source{Ours (monotonicity as in Definition~\ref{bp:p2s2:constantsB}; the only division, the cubic iso term, keeps its sup denominator).}
\end{lemma}

%----------------------------------------------------------------------------------------------------
\subsection{Engine bookkeeping}\label{bp:p2s2l1:sec-engine}

\begin{proposition}[The entry in the stage engine]\label{bp:p2s2l1:engine}
\begin{enumerate}[label=(\alph*)]
\item \emph{(entry by hand, then the engine)} Proposition~\ref{bp:p2s2l1:entryA} or~\ref{bp:p2s2l1:entryB} gives $K_1$ admissible on $T_{\rho_1}$, $\rho_1:=\rho-3\delta_0$,
  with error $\le$ the entry recursion value, ball $\beta-D_0$ and caps $m^{(0)}_i+\varrho^{(0)}_iw_i(\delta_0)X_0$ for every cap $i$ (frame, Gram inverse,
  $L^\top,N,N^\top$, $\abs T$, $\avg T^{-1}$, and the frequency in the iso form). The engine of Proposition~\ref{bp:p2s2:engine} runs on the
  \emph{original} schedule from stage $1$ (bites $\delta_s=\delta_0/2^s$, strips $r_{s+1}=r_s-3\delta_s$, limit strip $\rho-6\delta_0\ge\rho_\infty$): the entry is
  stage $0$ of that schedule, not a restart from $\rho_1$ with a recomputed capped bite $\min\{1,(\rho_1-\rho_\infty)/6\}$, which equals $\delta_0/2$ only when
  $\rho-\rho_\infty\le6$ (audit AUDIT-B10D-1, F5). In Lean the run is route A's stage-indexed engine (SHARP2-B9, \Lean{KAM.Sharp2.RouteA.SchemeA.run}),
  whose stage $0$ carries the entry constants; it gives the limit; its hypotheses \Lean{hnn}, \Lean{hcondT} (the step conditions are $\mathrm{coef}\cdot X\le\mathrm{bound}$ with $\mathrm{coef}\ge0$, so
  downward closed) and \Lean{hball} (the total displacement, entry plus engine, strictly below $\beta$) are part of the smallness predicate.
\item \emph{(route A in the route-B-normalised engine)} For route A put $X:=\varepsilon/(\gamma^4\delta^{4\varkappa})$, $\mu:=\gamma\delta^{\varkappa+1}X=\varepsilon/(\gamma^3\delta^{3\varkappa-1})$,
  $Q(\delta):=2^{3\varkappa-1}C_E(\delta)$ and the engine displacement constant $\gamma^2\delta^{2\varkappa-1}C_{\Delta K}(\delta)$; then the engine recursion
  $X_{s+1}=2^{\varkappa+1}Q_sX_s^2$ is $x_{s+1}=2^{4\varkappa}C_{E,s}x_s^2$ (audit F5). The run is indexed by the stage $s$, which fixes the bite $\delta_s$
  and the strip $r_s$ of the schedule; no bite function of the strip and no non-schedule bite occurs (route A's engine of SHARP2-B9, audit
  AUDIT-B10D-1, F4). The frame weights are $\gamma^2\delta^{2\varkappa-1}$ and the ball weight $\gamma^2\delta^{2\varkappa}$; in a body uniform in
  $\gamma\le\bar\gamma$ the cap sums are checked with $\bar\gamma^2$.
\item \emph{(closeness)} The ball sum is $D\le\gamma^2(x_0/X_\star)\tilde S(X_\star)$ with $\tilde S:=\sum_sC_{\Delta K,s}\delta_s^{2\varkappa}X_s$ (route A) and
  $D\le(X_0/X_\star)S(X_\star)$, $S$ the ball sum of (J3) (route B); hence the frozen closeness constants $C_2=\rho^{2\tau}\tilde S/(X_\star\delta_0^{4\tau})$ (route A) and
  $C_2=\rho^{2\tau}\sigma_{L^\top}S/(X_\star\delta_0^{2\tau+1})$ (route B), with no $\bar\gamma^2$ (audit F6).
\end{enumerate}
\source{Proposition~\ref{bp:p2s2:engine}; Corollary~\ref{bp:p2s2:hybrid} (the restart principle); audits AUDIT-SHARP-V2-fixed-1 (F2, F5, F6) and
AUDIT-SHARP-V2-iso-1 (F2); ours.}
\end{proposition}

\begin{proof}
\proofstep{Restart}\label{step:p2s2l1:en-restart}
The entry is one application of the step, at stage $0$ of the schedule; the next bite is $\delta_1=\delta_0/2$ by the definition of the schedule,
with no condition on $\rho-\rho_\infty$. The caps of $K_1$ are the initial values plus the entry increments, and the engine's (J3) budgets are what
remains of the caps; the displacements add.

\proofstep{Route A}\label{step:p2s2l1:en-A}
$\mu_{s+1}=\varepsilon_{s+1}/(\gamma^3(\delta_s/2)^{3\varkappa-1})\le2^{3\varkappa-1}C_E\varepsilon_s^2/(\gamma^7\delta_s^{7\varkappa-1})=Q_s\mu_s^2/(\gamma\delta_s^{\varkappa+1})$; the displacement
$C_{\Delta K}\varepsilon/(\gamma^2\delta^{2\varkappa})$ equals $(\gamma^2\delta^{2\varkappa-1}C_{\Delta K})\mu/(\gamma\delta^\varkappa)$; the strip loss is $3\delta$ ($\lambda=2+\eta$, $\eta=1$).

\proofstep{Closeness}\label{step:p2s2l1:en-cl}
The recursions are increasing in $X_0$ and every increment is $\ge$ linear in it, so $S(X_0)\le(X_0/X_\star)S(X_\star)$ for $X_0\le X_\star$; substitute
$x_0=e/(\gamma^4\delta_0^{4\tau})$, resp.\ $X_0=pe/(\gamma\delta_0^{\tau+1})$ with $p\le\sigma_{L^\top}/(\gamma\delta_0^\tau)$ when the normal term dominates (else $C_2$ uses $p$).
\end{proof}

%----------------------------------------------------------------------------------------------------
\subsection{The theorems}\label{bp:p2s2l1:sec-g1}

\begin{theorem}[G1 with the $\ell^1$ entry, route A]\label{bp:p2s2l1:g1A}
Let $(h,K,\omega,\gamma,e)$ satisfy the hypotheses of Theorem~\ref{bp:p0:g1} except the smallness, the refined data \ref{bp:p2s2:R1}--\ref{bp:p2s2:R3}
with initial margins for every cap, and the entry data \ref{bp:p2s2l1:L1}--\ref{bp:p2s2l1:L5} (\ref{bp:p2s2l1:L4} optional). Let $X_\star>0$ satisfy the route-A
smallness predicate: the entry conditions at $X_\star$, the engine conditions (J1)--(J3) of Proposition~\ref{bp:p2s2l1:engine}(a), (b) from the entry
value $2^{4\varkappa}C_E^{\mathrm{entry}}X_\star^2$, with every cap budget reduced by the entry increment, and the ball condition. If
$e/(\gamma^4\delta_0^{4\tau})\le X_\star$, i.e.\ $C_1e/(\gamma^4\rho^{4\tau})\le1$ with $C_1:=\rho^{4\tau}/(X_\star\delta_0^{4\tau})$, then the conclusion of Theorem~\ref{bp:p0:g1}
holds with $\abs{K_\infty-K}\le C_2e/(\gamma^2\rho^{2\tau})$, $C_2$ as in Proposition~\ref{bp:p2s2l1:engine}(c).
\source{Theorem~\ref{bp:p2s2:g1} (route A); Propositions~\ref{bp:p2s2l1:entryA}, \ref{bp:p2s2l1:engine}; ours.}
\Lean{KAM.Sharp2.L1.g1\_v2A\_l1} (with the v1 and sup branches, SHARP2-B10C)
\end{theorem}

\begin{proof}
\proofstep{Entry, engine, limit}\label{step:p2s2l1:g1A}
At stage $0$ the error is $\abs E\le e$ and the entry data hold; Proposition~\ref{bp:p2s2l1:entryA} gives $K_1$; Proposition~\ref{bp:p2s2l1:engine}(a), (b) and the
limit steps of Proposition~\ref{bp:p2:limit} (as in Step~\ref{step:p2s2:g1-a}) give $K_\infty$, its invariance, margins and closeness. Monotonicity in $x_0$
(Lemma~\ref{bp:p2s2l1:never-worse} and Corollary~\ref{bp:p2s2:frozen-form}) gives the form with $X_\star$.
\end{proof}

\begin{theorem}[G1 with the $\ell^1$ entry, route B]\label{bp:p2s2l1:g1B}
The same with Proposition~\ref{bp:p2s2l1:entryB}, the route-B normalisation $X_0=pe/(\gamma\delta_0^{\tau+1})$, $C_1:=\rho^{4\tau}\bar\gamma^2\sigma_{L^\top}/(X_\star\delta_0^{2\tau+1})$ when
$p=\sigma_{L^\top}/(\gamma\delta_0^\tau)$ (the regime of every package; otherwise the explicit $p$), and the closeness of Proposition~\ref{bp:p2s2l1:engine}(c).
\source{Theorem~\ref{bp:p2s2:g1} (route B); Proposition~\ref{bp:p2s2l1:entryB}; ours.}
\end{theorem}

\begin{corollary}[Frozen form and the minimum with v2]\label{bp:p2s2l1:frozen}
For tori with certified entry data, the smallness of the v2-with-entry theorem is the disjunction of route A with the entry, route B with the
entry and the routes of Theorem~\ref{bp:p2s2:g1} (v1 and the sup routes, each with its own certificate); by Lemma~\ref{bp:p2s2l1:never-worse} the entry
routes are never worse than the sup routes for the same data. The final constant $C_1$ is the minimum of the branches, and the closeness constants
$C_2$ (and $C_3$ in the iso form) the maximum, proved by cases on the branch that applies (audits AUDIT-SHARP-V2-fixed-1 F6, AUDIT-B10D-1 F6). Tori
without entry data are covered by Theorem~\ref{bp:p2s2:g1} itself.
\end{corollary}

%----------------------------------------------------------------------------------------------------
\subsection{Numbers, deviations and register}\label{bp:p2s2l1:sec-register}

\begin{remark}[Register]\label{bp:p2s2l1:register}
\begin{enumerate}[label=(\roman*)]
\item \emph{PlasmaLean's example} (\texttt{dominance.py --plasmalean}; its tolerance $10^{-18.24}$; $\omega=(\iota,1)$, $\gamma=0.0687$, $\rho=0.02$; its declared
  bounds read as sup and $\ell^1$ bounds, transposes at the conservative defaults; sup stages with $c_{\mathrm{freq}}$ computed at each bite,
  $0.00754$ to $0.00752$; $c^\infty_0=0.00126$):
  \begin{center}\small
  \begin{tabular}{lrrrr}
  & \multicolumn{2}{c}{(a) $\ell^1$ vs $\ell^1$} & \multicolumn{2}{c}{(b) sup-equivalent, $\ell^1/\text{sup}=3.6$}\\
  variant & $\log_{10}e^*$ & margin & $\log_{10}e^*$ & margin\\\hline
  A7 (sup, v2) & $-17.27$ & $+0.97$ & $-17.27$ & $+0.97$\\
  A7 + entry & $-15.91$ & $+2.33$ & $-16.44$ & $+1.80$\\
  A7 + entry + $\ell^1$ frame increment & $\mathbf{-15.40}$ & $\mathbf{+2.84}$ & $\mathbf{-15.93}$ & $\mathbf{+2.31}$\\
  B6 (sup, torsion cap) & $-18.10$ & $+0.14$ & $-18.10$ & $+0.14$\\
  B6 + entry & $-17.12$ & $+1.12$ & $-17.65$ & $+0.59$
  \end{tabular}
  \end{center}
  These margins are \emph{conditional comparisons of scalar error tolerances at assumed data} ($\gamma=0.0687$, $\rho\approx0.02$, the clock
  frequency, and guessed canonical frame and torsion norms, with the declared bounds read as both sup and $\ell^1$ bounds, $\nu_K=\sigma_K$,
  $\nu_N=\sigma_N$; a sup bound alone does not give that reading). Reading (a) compares our threshold on $\norm E_{1,\rho}$ with PlasmaLean's
  tolerance on the $\ell^1$ error-norm category, not on an identified canonical residual; reading (b) assumes the ratio $\ell^1/\text{sup}=3.6$ of
  the R1 lower torus (CERTIFIER-SHARPNESS \S1.2; not a measurement or enclosure for PlasmaLean's torus) and reruns the mirror with $a_N=a_L=3.6$, so
  it differs from (a) by $0.53$ orders, not by $\log_{10}3.6$. With the ``R1-like'' norms the margins are $+6.18$ and $+5.73$. PlasmaLean's field is
  non-canonical ($1\frac12$ degrees of freedom) and no canonical chart or transfer of its residual and norms is supplied, so the owner's
  application-level dominance requirement is not established by these rows (audit AUDIT-B10D-1, F7).
\item \emph{R1} (\texttt{thresholds\_v2.py --l1}; measured $\ell^1$ data $\times1.05$, heuristic until enclosed): route A $-11.18\to-10.42$ (entry)
  $\to-10.13$ (with the frame increment), iso $-11.37\to-10.60\to-10.32$; route B $-12.25\to-12.16$, iso $-12.39\to-12.31$. Sensitivity to the
  ratio $a=a_N=a_L$ (fixed, route A, all levers): $a=1$: $-9.67$; $2$: $-9.91$; $3.6$: $-10.11$; $8$: $-10.29$. Binding: the $\avg T^{-1}$ cap sum
  (route A), the tail decay (route B).
\item \emph{Comparison with the emulation of the design document} (\S10.1, ``$\ell^1$ (L9)'' rows: $c_\infty$ at every stage, $E$ charged $\times3.6$): $-16.24$
  for PlasmaLean (cA-like). The entry alone gives $-16.44$ in the same reading, without any conversion; the $\ell^1$ frame increment, absent from the
  emulation, gives $-15.93$.
\item \emph{Deviations} (ours): the entry step itself (no source runs one Newton step in $\ell^1$ and the rest in sup); the per-ingredient minimum; the
  $\ell^1$ frame increment; the defect identity used constructively (Lemma~\ref{bp:p2s2l1:defect}); the full-bite average in $\ell^1$ at real points.
\item \emph{Inherited open points.} Audit AUDIT-SHARP-V2-fixed-1 F8 (the torsion of the new torus and the stage transition) applies unchanged to the
  $\mathcal T_c$, $\mathcal T_F$ terms of $\bar K$ in the entry step, which keeps the sup treatment of Part~\ref{bp:p2s2}; AUDIT-SHARP-V2-iso-1 F5 (moving-frequency
  factors) applies unchanged to the iso entry. Their repairs change the sup constants, and the entry inherits them through Definitions~\ref{bp:p2s2l1:constA},
  \ref{bp:p2s2l1:constB}.
\item \emph{Suspected errata.} None in a pinned source. Our own: the design document \S10.2 says the hybrid needs ``a parallel norm layer for the step''
  (1.5--2k lines); the entry design needs no $\ell^1$ step layer beyond stage $0$.
\end{enumerate}
\end{remark}
